%% =============================================================================
%% section2.tex — Section 2: Motivating Data and Preliminary Analysis
%% =============================================================================

\section{Motivating Data and Preliminary Analysis}
\label{sec:data}

This section introduces the data and documents the empirical features
that jointly determine the model developed in \cref{sec:model}.  We
organize the presentation around four distributional
characteristics---structural zeros, bounded support, overdispersion, and
complex survey design---each of which maps directly to a component of the
hierarchical hurdle beta-binomial specification.  A preliminary analysis
then establishes the reversal phenomenon and the cross-state
heterogeneity that motivate the hierarchical structure and cross-margin
covariance.

%% ---------------------------------------------------------------------------
%% 2.1 The NSECE Survey
%% ---------------------------------------------------------------------------

\subsection{The NSECE survey}\label{sec:nsece}

The 2019 National Survey of Early Care and Education
\citep[NSECE;][]{NSECEProjectTeam2022} is a nationally representative
survey of early care and education providers conducted by NORC at the
University of Chicago under contract with the Office of Planning,
Research, and Evaluation (OPRE) within the Administration for Children
and Families.  Data collection took place from November 2018 through
July 2019, with $6{,}917$ center-based providers completing the survey.
The center-based provider component employs a stratified
multistage cluster design: 30 strata defined by state groupings, 415
primary sampling units, and $6{,}809$ centers with nonmissing survey
responses.
Of these, 24 report a total enrollment of $n_i = 0$ and are excluded
because the beta-binomial kernel requires a positive denominator,
yielding a final analytic sample of $N = 6{,}785$ providers across all
$S = 51$ states (50 states plus the District of Columbia).  Full details of the sampling frame, weight
construction, and informativeness testing appear in Supplementary
Material~D.

For each provider $i$, the survey records total enrollment $n_i$ (ages
0--5) and infant/toddler (IT) enrollment $Y_i \in \{0, 1, \ldots, n_i\}$
(ages 0--2).  The outcome pair $(Y_i, n_i)$ is the natural unit of
analysis.  Provider-level covariates include the community poverty rate,
urbanicity, and racial/ethnic composition of the provider's Census
tract~\citep{DobbinsEtAl2016, MalikEtAl2018}.  State-level policy
variables capture three dimensions of Child Care and Development Fund
(CCDF) subsidy policy---market rate percentile, tiered reimbursement, and
infant/toddler rate add-on---and enter the model as cross-level
moderators in \cref{sec:application}.  \Cref{tab:data-summary}
summarizes the key descriptive statistics.

\begin{table}[t]
\centering
\caption{Data summary: NSECE 2019 center-based providers.
  Panel~A describes the outcome variables; Panel~B the provider-level
  covariates (continuous variables shown before standardization, as
  means with standard deviations and ranges); Panel~C the survey design
  features, including the Kish design effect
  $\DEFF_{\mathrm{Kish}} = 1 + \mathrm{CV}^2(w)$.}\label{tab:data-summary}
\smallskip
\begin{tabular}{@{}ll@{}}
\toprule
Variable & Value \\
\midrule
\multicolumn{2}{@{}l}{\textit{Panel A: Outcome}} \\[2pt]
Total providers ($N$) & 6{,}785 \\
IT-serving providers ($z_i = 1$) & 4{,}392 (64.7\%) \\
Non-servers ($z_i = 0$) & 2{,}393 (35.3\%) \\
Total enrollment ($n_i$): Mean / Median / Range & 60.5 / 48 / 1--378 \\
IT enrollment ($Y_i \mid z_i = 1$): Mean / Median & 29.4 / 24 \\
IT share ($Y_i/n_i \mid z_i = 1$): Mean (SD) & 0.478 (0.214) \\[4pt]
\multicolumn{2}{@{}l}{\textit{Panel B: Covariates --- Mean (SD) [Range]}} \\[2pt]
Community poverty rate (\%) & 17.3 (8.3) [1.9, 54.2] \\
Urban (\%) & 93.3 \\
Community \% Black & 17.2 (20.5) [0.0, 96.5] \\
Community \% Hispanic & 24.2 (22.6) [0.0, 98.4] \\[4pt]
\multicolumn{2}{@{}l}{\textit{Panel C: Survey design}} \\[2pt]
States ($S$) & 51 (50 + DC) \\
Strata / PSUs & 30 / 415 \\
Sampling weight range & 1--462 \\
Weight CV & 1.66 \\
Kish effective sample size & 1{,}803 \\
Kish $\DEFF$ & 3.76 \\
\bottomrule
\end{tabular}
\end{table}

%% ---------------------------------------------------------------------------
%% 2.2 Four Data Features
%% ---------------------------------------------------------------------------

\subsection{Four data features and their modeling implications}\label{sec:features}

We now examine four distributional characteristics of the NSECE data
that individually and jointly determine the model specification of
\cref{sec:model}.  \Cref{fig:it-distribution} displays the marginal
distribution of the infant/toddler enrollment share, which exhibits all
four features simultaneously.

\begin{figure}[t]
  \centering
  \includegraphics[width=\textwidth]{F1_it_distribution.pdf}
  \caption{Distribution of the infant/toddler enrollment share $Y_i/n_i$
    across $N = 6{,}785$ center-based providers (NSECE 2019). The spike at
    zero (35.3\% of centers) represents providers that do not serve any
    children under age 3. Among IT-serving providers ($z_i = 1$, $N_{\mathrm{pos}} = 4{,}392$),
    the distribution is roughly unimodal with a mean of 0.478 and substantial
    dispersion (SD $= 0.214$), consistent with overdispersion approximately
    twelve times beyond the binomial---three features that jointly motivate
    the hurdle beta-binomial model.}\label{fig:it-distribution}
\end{figure}

\paragraph{Structural zeros and the hurdle specification.}
Of the 6,785 providers, 2,393 (35.3\%) report zero infant/toddler
enrollment---a rate far too large for sampling variability and reflecting
a deliberate operational decision not to serve children under age
three~\citep{AdamsRohacek2010}.  The hurdle
specification~\citep{Mullahy1986} treats participation and intensity as
distinct processes governed by separate parameters, as developed in
\cref{sec:likelihood}.

\paragraph{Bounded support and the beta-binomial kernel.}
Among servers ($z_i = 1$), infant/toddler enrollment $Y_i$ is bounded
above by total enrollment $n_i$, so the share $Y_i/n_i$ lies in
$(0,1]$.  The beta-binomial distribution respects this
$\{0,1,\ldots,n_i\}$ support and incorporates extra-binomial variation
through the concentration parameter
$\kappa$~\citep{Prentice1986, WagnerEtAl2015}; see
\cref{sec:likelihood}.

\paragraph{Overdispersion and the concentration parameter $\kappa$.}
The observed variance of $Y_i/n_i$ among participants exceeds the
binomial prediction by a factor of approximately 12 at the pooled level
(marginal overdispersion, before accounting for covariates),
and all 51 states individually exhibit overdispersion.  After conditioning
on the model covariates and hierarchical structure, the estimated
concentration $\hat{\kappa} \approx 7$ implies a conditional overdispersion
factor of approximately 7.  The beta-binomial
nests the binomial as $\kappa \to \infty$, modeling this overdispersion
structurally rather than correcting for it post hoc~\citep{Prentice1986}.

\paragraph{Complex survey design and the pseudo-posterior.}
Sampling weights range from 1 to 462, yielding a Kish effective sample
size of $\ESS_{\mathrm{Kish}} = 1{,}803$ and a design effect of
$\DEFF_{\mathrm{Kish}} = 3.76$; a formal test confirms that the design
is informative for the intensive margin~\citep{Pfeffermann1993}.
Ignoring the design would yield credible intervals with empirical
coverage well below nominal; \cref{sec:survey} develops the
pseudo-posterior and sandwich correction.

Each feature alone eliminates a class of candidate models: structural
zeros eliminate single-equation specifications, bounded support
eliminates Poisson and negative-binomial kernels, universal
overdispersion eliminates the binomial, and the complex survey design
eliminates naive Bayesian posteriors.  The \emph{intersection} of all
four constraints narrows the viable specification to a hurdle model with
a beta-binomial kernel, estimated via a pseudo-posterior with sandwich
correction.

%% ---------------------------------------------------------------------------
%% 2.3 Preliminary Evidence for the Poverty Reversal
%% ---------------------------------------------------------------------------

\subsection{Preliminary evidence for the poverty reversal}\label{sec:prelim}

The four data features determine the model's distributional
specification.  The model's \emph{hierarchical} structure---state-varying
coefficients with cross-margin covariance---is motivated by a fifth
empirical pattern: the reversal varies systematically across
states.

\paragraph{Pooled estimates.}
A pooled logistic regression of infant/toddler participation ($z_i$) on
the standardized community poverty rate yields an extensive-margin
coefficient of $\hat{\alpha}_{\mathrm{pov}} = -0.020$
($p < 10^{-10}$), confirming that providers in higher-poverty
communities are significantly less likely to serve infants.  A pooled OLS
regression of the enrollment share ($Y_i/n_i$ among servers) on poverty
yields an intensive-margin coefficient of
$\hat{\beta}_{\mathrm{pov}} = +0.003$ ($p < 10^{-14}$), establishing
that, conditional on serving, providers in poorer communities devote a
larger share of enrollment to infants.  The opposing signs constitute the
reversal in its simplest form: poverty \emph{reduces}
participation but \emph{increases} intensity.

The substantive interpretation is noteworthy.  Poverty erects a barrier to
entry: the elevated costs of infant care discourage participation
precisely where demand is greatest~\citep{DobbinsEtAl2016,
MalikEtAl2018}.  Yet among centers that clear this barrier, those in
disadvantaged communities tilt their enrollment mix toward the youngest
age group, reflecting either market demand in areas with fewer
alternatives or mission-driven service provision.  Any model that
collapses participation and intensity into a single equation would
estimate a muted net effect, obscuring the two countervailing forces at
work.

\paragraph{State heterogeneity.}
When these regressions are estimated separately by state, 23 of 51
states (45\%) display the classic reversal pattern
($\hat{\alpha}_{\mathrm{pov},s} < 0$ and
$\hat{\beta}_{\mathrm{pov},s} > 0$).  The remaining states show varying
combinations: some exhibit negative effects on both margins, while
others show positive effects on both.  The cross-state correlation
between the two margin-specific poverty coefficients is $r = 0.021$---near
zero---indicating that the extensive and intensive mechanisms operate
largely independently.  \Cref{fig:poverty-raw} displays the aggregate reversal
in the raw data; the state-level bivariate scatter appears
in~\cref{fig:cross-margin}.  The near-zero cross-margin correlation implies
that knowing a state has a strong participation barrier provides
essentially no information about its intensity response, a finding that
motivates the full covariance structure developed in
\cref{sec:model}~\citep{SippleEtAl2020}.

\begin{figure}[t]
  \centering
  \includegraphics[width=\textwidth]{F2_poverty_reversal_raw.pdf}
  \caption{Preliminary evidence for the poverty reversal in the raw NSECE
    2019 data.  \emph{Left panel:} IT participation rate (fraction of
    centers serving any infants) by community poverty decile; a loess
    smoother with 95\% confidence band confirms the declining trend.
    \emph{Right panel:} mean IT enrollment share ($Y_i/n_i$) among servers
    by poverty decile, showing the opposing positive trend.  The two panels
    together display the reversal in its simplest form: participation falls
    with poverty, but intensity rises.}\label{fig:poverty-raw}
\end{figure}

\paragraph{Policy correlations.}
Among the state-level policy variables, tiered reimbursement shows the
strongest association with the poverty--participation relationship:
states with tiered reimbursement tend to have less negative
extensive-margin poverty coefficients, consistent with the hypothesis
that quality-linked payment incentives offset the cost disadvantage of
serving infants in disadvantaged communities.  These preliminary
correlations motivate the cross-level moderation analysis in
\cref{sec:application}.

In sum, the data require a model that simultaneously (i)~separates
structural zeros from positive counts via a hurdle, (ii)~respects the
bounded support via a beta-binomial kernel, (iii)~accommodates
overdispersion through the concentration parameter, (iv)~incorporates
survey weights through a pseudo-posterior with sandwich correction, and
(v)~captures cross-state heterogeneity and cross-margin dependence
through a hierarchical prior with full covariance.  \Cref{sec:model}
develops this model.
